\subsection*{Data Expression}

So far, you have only used simple data expressions, either a defined \textbf{spanned data} or a \textbf{storage location}. The \textbf{storage location} is used solely for the \textbf{destination buffer} declaration, as it requires the Croqtile compiler to allocate memory for it. The following code provides an example:

\begin{lstlisting}[language=croqtile,caption={Data Expression}]
global f32 [32, 16] input;
shared f32 [32, 16] output;
dma.copy input => output;  // Copy to a user-declared buffer
dma.copy input => shared;  // Same, but to a compiler-allocated buffer
\end{lstlisting}

In this example, the destination of the second DMA operation is specified as the storage location \texttt{shared}. This requires the compiler to allocate storage accordingly, making it equivalent to the first DMA operation. In practice, you recommend programmers use storage locations as destinations, as Croqtile can deduce the destination shape from the \textbf{DMA statement}.

However, the use of simple \textbf{data expressions} is rare in practice. In the next section, this section introduces the \textbf{ChunkAt} expression to demonstrate how to implement tiling in data expressions.
